\newpage
\chapter{TINJAUAN PUSTAKA} \label{Bab II}

\section{Tinjauan Pustaka} \label{II.Tinjauan}
Bagian ini membahas penelitian terdahulu yang relevan dengan topik \textit{fine-tuning} Whisper dan Automatic Speech Recognition (ASR) untuk bahasa \textit{low-resource}. Pembahasan disusun secara sistematis dengan: (i) masalah yang diangkat, (ii) metode yang digunakan, (iii) hasil/pengujian, (iv) analisis kelebihan/keterbatasan, serta (v) perbandingan dengan penelitian yang dikerjakan pada skripsi ini (Whisper Small untuk Bahasa Minangkabau). 

\subsection*{Literatur Terkait (2022--2025)}
\begin{enumerate}[noitemsep,topsep=2pt]

\item \textbf{Robust Speech Recognition via Large-Scale Weak Supervision} \cite{radford2023whisper}. Radford, Kim, Xu, Brockman, McLeavey, Sutskever (2023).

\textbf{Masalah:} Bagaimana membangun model ASR multibahasa yang \textit{robust} dan mampu \textit{zero-shot} pada beragam kondisi tanpa \textit{fine-tuning} dataset spesifik.

\textbf{Metode:} Pra-latih \textit{Whisper} pada $\sim$680{,}000 jam data multibahasa/multitugas \textit{weak supervision}.

\textbf{Hasil/Pengujian:} Kinerja kuat di berbagai benchmark ASR/terjemahan wicara dalam skenario \textit{zero-shot}.

\textbf{Analisis:} \textbf{(+)} Fondasi kokoh untuk adaptasi ke bahasa lokal. \textbf{(--)} Bahasa/dialek yang under-represented masih memunculkan kesalahan.

\textbf{Perbandingan:} Menjadi dasar untuk \textbf{adaptasi} ke Minangkabau melalui \textit{fine-tuning} (Whisper Small) pada penelitian ini.

\item \textbf{Exploration of Whisper Fine-Tuning Strategies for Low-Resource ASR} \cite{liu2024exploration}. Liu, Yang, Qu (2024).

\textbf{Masalah:} Strategi \textit{fine-tuning} Whisper yang paling efektif untuk skenario data terbatas belum dikaji komprehensif.

\textbf{Metode:} Studi komparatif beberapa strategi (mis.\ \textit{full fine-tuning}, \textit{partial/freezing}, \textit{re-init layer atas}, \textit{adapters}, PEFT/LoRA) pada banyak bahasa \textit{low-resource}.

\textbf{Hasil/Pengujian:} Semua strategi menurunkan WER; terdapat \textit{trade-off} akurasi vs efisiensi (VRAM/waktu latih).

\textbf{Analisis:} \textbf{(+)} Panduan praktis memilih strategi berdasar keterbatasan data/komputasi. \textbf{(--)} Tidak mendalami satu bahasa/dialek spesifik secara longitudinal.

\textbf{Perbandingan:} Skripsi ini akan membandingkan \textbf{full FT} vs \textbf{LoRA} pada Whisper Small untuk Minangkabau.

\item \textbf{Fine-tuning Whisper on Low-Resource Languages for Real-World Applications} \cite{timmel2024finetune}. Timmel, Paonessa, Kakooee, Vogel, Perruchoud (2024).

\textbf{Masalah:} Kekurangan korpus \textit{long-form} membuat model kehilangan kemampuan segmentasi/\textit{timestamping} saat hanya dilatih pada kalimat pendek.

\textbf{Metode:} Menyusun ulang data kalimat pendek menjadi korpus audio \textit{long-form} sintetis (kasus Swiss German), lalu \textit{fine-tune} Whisper.

\textbf{Hasil/Pengujian:} Peningkatan signifikan pada aplikasi dunia nyata (percakapan panjang); kemampuan segmentasi tetap terjaga.

\textbf{Analisis:} \textbf{(+)} Solusi praktis ketika \textit{long-form} asli minim. \textbf{(--)} Diverifikasi pada Swiss German; butuh validasi lintas bahasa daerah lain.

\textbf{Perbandingan:} Relevan untuk Minangkabau bila data \textit{long-form} kurang; dapat meniru strategi penyusunan \textit{long-form} sintetis.

\item \textbf{Towards Efficiently Whisper Fine-tuning with Monotonic Alignments} \cite{zhuang2025monotonic}. Zhuang, Wei, Fang, Cheng, Wang, Xiao (2025).

\textbf{Masalah:} \textit{Fine-tuning} standar tidak mengoptimalkan keselarasan monotonic audio$\to$teks secara eksplisit.

\textbf{Metode:} Tambah \textit{soft monotonic alignment loss} saat \textit{fine-tuning} Whisper, tanpa mengubah arsitektur.

\textbf{Hasil/Pengujian:} Penurunan WER konsisten pada beberapa tugas ASR multibahasa; perbaikan juga pada \textit{speech translation}.

\textbf{Analisis:} \textbf{(+)} Peningkatan akurasi dengan overhead kecil. \textbf{(--)} Perlu setel \textit{loss weight}; dampak di \textit{very low-resource} perlu studi lanjut.

\textbf{Perbandingan:} Menjadi opsi \textbf{pengembangan} jika performa baseline Minangkabau belum memadai.

\item \textbf{Indonesian Automatic Speech Recognition with XLSR-53} \cite{arisaputra2022xlsr53}. Arisaputra, Zahra (2022).

\textbf{Masalah:} Mencapai WER kompetitif Bahasa Indonesia dengan data terbatas.

\textbf{Metode:} \textit{Fine-tuning} XLSR-53 ($\sim$24 jam) + \textit{language model} (KenLM 4-gram) pada decoding.

\textbf{Hasil/Pengujian:} WER $\sim$20\% $\to$ $\sim$12\% dengan LM pada \textit{test set} publik.

\textbf{Analisis:} \textbf{(+)} Bukti kuat \textit{transfer learning} multibahasa. \textbf{(--)} Fokus BI baku; tidak membahas dialek/daerah.

\textbf{Perbandingan:} Menjadi pembanding non-Whisper; mendukung hipotesis bahwa model pra-latih multibahasa + FT efektif untuk bahasa Indonesia/daerah.

\item \textbf{Breaking the Transcription Bottleneck: Fine-tuning ASR Models for Extremely Low-Resource Fieldwork Languages} \cite{liang2025fieldmatters}. Liang, Levow (2025).

\textbf{Masalah:} Transkripsi bahasa yang sangat minim data (bahkan $<$1 jam) dalam konteks linguistik lapangan.

\textbf{Metode:} Bandingkan \textbf{MMS} vs \textbf{XLS-R} di 5 bahasa sangat \textit{low-resource} pada berbagai skala data.

\textbf{Hasil/Pengujian:} MMS unggul pada $<$1 jam; paritas dengan XLS-R saat data bertambah ($\gtrsim$1 jam).

\textbf{Analisis:} \textbf{(+)} Panduan pemilihan model sesuai ketersediaan data. \textbf{(--)} Spesifik MMS/XLS-R; bukan Whisper.

\textbf{Perbandingan:} Memberi konteks batas bawah data; skripsi ini tetap memfokuskan Whisper Small + FT untuk Minangkabau.

\end{enumerate}

\subsection*{Rangkuman \& Peluang Pengembangan}
Literatur terbaru menegaskan bahwa: (i) pra-latih skala besar \cite{radford2023whisper} memberi fondasi kuat, (ii) pemilihan strategi FT harus menimbang akurasi vs efisiensi \cite{liu2024exploration}, (iii) rekayasa korpus \textit{long-form} membantu aplikasi dunia nyata \cite{timmel2024finetune}, (iv) modifikasi objektif (\textit{alignment loss}) dapat meningkatkan akurasi tanpa mengubah arsitektur \cite{zhuang2025monotonic}, dan (v) pembelajaran multibahasa dapat ditransfer ke bahasa lokal/daerah \cite{arisaputra2022xlsr53}. Untuk penelitian ini, peluang pengembangan mencakup: (a) komparasi \textbf{full FT vs LoRA} pada Whisper Small (efisiensi VRAM), (b) penyusunan \textbf{korpus long-form sintetis} Minangkabau jika data panjang minim, (c) eksplorasi \textbf{alignment-based loss} dan/atau \textbf{language model} pada decoding untuk penurunan WER lanjutan.

{\small
\setlength{\tabcolsep}{4pt}% padding kolom lebih rapat
\begin{tabularx}{\textwidth}{|>{\centering\arraybackslash}m{0.045\textwidth}|Y|Y|Y|Y|}
\caption{Ringkasan Literasi Penelitian Terdahulu (2022--2025)}
\label{table:2.literasi}\\
\hline
\textbf{No.} & \textbf{Judul} & \textbf{Masalah} & \textbf{Metode} & \textbf{Hasil (ringkas)}\\
\hline
\endfirsthead
\hline
\textbf{No.} & \textbf{Judul} & \textbf{Masalah} & \textbf{Metode} & \textbf{Hasil (ringkas)}\\
\hline
\endhead
1 &
\textit{Robust Speech Recognition via Large-Scale Weak Supervision} \cite{radford2023whisper} &
Generalisasi \textit{zero-shot} &
Pra-latih $\sim$680k jam multibahasa/
multitugas &
Robust di banyak benchmark; dasar adaptasi lokal \\
\hline
2 &
\textit{Exploration of Whisper Fine-Tuning Strategies for Low-Resource ASR} \cite{liu2024exploration} &
Strategi FT untuk \textit{low-resource} &
Bandingkan full/partial
/adapters/LoRA &
Semua turunkan WER; ada \textit{trade-off} akurasi vs efisiensi \\
\hline
3 &
\textit{Fine-tuning Whisper on Low-Resource Languages for Real-World Applications} \cite{timmel2024finetune} &
Minim \textit{long-form} $\to$ hilang segmentasi &
Susun kalimat $\to$ \textit{long-form} sintetis + FT &
SOTA Swiss German; segmentasi/
\textit{timestamping} tetap baik \\
\hline
4 &
\textit{Towards Efficiently Whisper Fine-tuning with Monotonic Alignments} \cite{zhuang2025monotonic} &
Alignment tak dioptimalkan &
Tambah \textit{monotonic alignment loss} &
WER turun konsisten tanpa ubah arsitektur \\
\hline
5 &
\textit{Indonesian Automatic Speech Recognition with XLSR-53} \cite{arisaputra2022xlsr53} &
ASR Indonesia data terbatas &
FT XLSR-53 + LM 4-gram &
WER $\sim$20\% $\to$ $\sim$12\% (dengan LM) \\
\hline
6 &
\textit{Breaking the Transcription Bottleneck} \cite{liang2025fieldmatters} &
Sangat minim data ({$<\!$}1 jam) &
Benchmark MMS vs.\ XLS\textendash R &
MMS unggul {$<\!$}1 jam; paritas saat data bertambah \\
\hline
\end{tabularx}
}

\section{Dasar Teori} \label{II.Teori}
Bagian ini membahas konsep inti yang mendasari pengembangan \textit{Automatic Speech Recognition} (ASR) modern (2020–2025), dengan fokus pada adaptasi Whisper untuk Bahasa Minangkabau. Uraian dimulai dari formulasi umum ASR, representasi fitur akustik, arsitektur Transformer modern, strategi \textit{fine-tuning} pada skenario \textit{low-resource}, decoding dengan bantuan model bahasa, hingga metrik evaluasi.

% ---------- 2.1 Formulasi umum ASR ----------
\subsection{Formulasi Umum ASR}
ASR memetakan fitur akustik $X$ menjadi urutan token teks $y=(y_1,\dots,y_T)$ dengan memaksimalkan probabilitas posterior.

\begin{eqfloat}[H]
\caption{Objektif umum ASR: pemilihan urutan hipotesis terbaik}
\label{eq:asr-objective}
\begin{equation}
\hat{y} \;=\; \arg\max_{y}\; P_\theta\!\left(y\,\middle|\,X\right).
\end{equation}
\end{eqfloat}

Persamaan~\ref{eq:asr-objective} adalah target inferensi untuk model \textit{encoder–decoder} modern seperti Conformer dan Whisper \cite{gulati2020conformer,radford2023whisper}.

% ---------- 2.2 Representasi akustik ----------
\subsection{Representasi Akustik (Log–Mel Spectrogram)}
Sebagian besar sistem terkini (termasuk Whisper) mengonsumsi \textit{log–Mel spectrogram}.

\begin{eqfloat}[H]
\caption{Pipeline komputasi log–Mel spectrogram}
\label{eq:logmel}
\begin{equation}
\begin{aligned}
S(k,m) &= \mathrm{STFT}\{s[n]\},\\
E(k,m) &= |S(k,m)|^2,\\
M(f,m) &= \mathrm{MelFilterbank}\!\left(E(k,m)\right),\\
X(f,m) &= \log\!\left(M(f,m)+\epsilon\right).
\end{aligned}
\end{equation}
\end{eqfloat}

Persamaan~\ref{eq:logmel} memetakan sinyal waktu $s[n]$ ke fitur $X$ yang lebih stabil terhadap variasi energi dan sesuai untuk encoder Transformer \cite{radford2023whisper}.

% ---------- 2.3 Arsitektur Transformer modern ----------
\subsection{Arsitektur Transformer Modern untuk ASR (2020–2025)}
\textbf{Conformer} menggabungkan konvolusi (lokal) dan \textit{self-attention} (global), efektif untuk ketergantungan pendek–panjang \cite{gulati2020conformer}.  
Pada jalur pra-latih \textit{self-supervised}, \textbf{wav2vec\,2.0} \cite{baevski2020wav2vec2}, \textbf{XLS-R} \cite{conneau2021xlsr}, dan \textbf{WavLM} \cite{chen2022wavlm} menyediakan representasi akustik lintas-bahasa yang kuat untuk diadaptasi dengan transkrip terbatas.

\begin{eqfloat}[H]
\caption{Skor perhatian (scaled dot-product) pada Transformer}
\label{eq:attn}
\begin{equation}
\mathrm{Attention}(Q,K,V)\;=\;\mathrm{softmax}\!\left(\frac{QK^\top}{\sqrt{d_k}}\right)\!V .
\end{equation}
\end{eqfloat}

% ---------- 2.4 Model Whisper ----------
\subsection{Model Whisper}
Whisper adalah \textit{encoder–decoder} Transformer multibahasa yang dipra-latih $\sim$680k jam \cite{radford2023whisper}. Pada adaptasi, digunakan kerugian \textit{cross-entropy} autoregresif:

\begin{eqfloat}[H]
\caption{Fungsi kerugian \textit{cross-entropy} untuk decoding autoregresif}
\label{eq:ce}
\begin{equation}
\mathcal{L}_{\mathrm{CE}} \;=\; -\sum_{t=1}^{T}\log P_\theta\!\left(y_t \,\middle|\, y_{<t}, X\right).
\end{equation}
\end{eqfloat}

Whisper memakai BPE multibahasa dan token kontrol (bahasa, modus tugas) untuk mengondisikan output \cite{radford2023whisper}.

% ---------- 2.5 Fine-tuning low-resource ----------
\subsection{Strategi Fine-Tuning pada Skenario Low-Resource}
\paragraph{Fine-tuning penuh.}
Seluruh parameter dioptimasi dengan \textit{learning rate} kecil; akurat namun berbiaya VRAM/waktu lebih tinggi.

\paragraph{Parameter-Efficient Fine-Tuning (LoRA).}
LoRA menambah pembaruan ber-rank rendah pada proyeksi linear (mis.\ di modul atensi) tanpa memodifikasi bobot pra-latih utama.

\begin{eqfloat}[H]
\caption{Formulasi LoRA untuk pembaruan ber-rank rendah}
\label{eq:lora}
\begin{equation}
W' \;=\; W + \Delta W,\qquad 
\Delta W \;=\; B A,\quad
A\!\in\!\mathbb{R}^{r\times k},\; B\!\in\!\mathbb{R}^{d\times r},\; r\ll\min(d,k).
\end{equation}
\end{eqfloat}

Studi 2024 menunjukkan semua strategi (penuh, \textit{partial freeze}, adapter, LoRA) menurunkan WER; pemilihan bergantung data dan sumber daya \cite{liu2024exploration}.

\paragraph{Kendala Alignment Monotonik.}
Untuk ASR, keselarasan audio$\rightarrow$teks cenderung monotonik; menambahkan penalti alignment saat \textit{fine-tuning} terbukti menurunkan WER (2025) \cite{zhuang2025monotonic}.

\begin{eqfloat}[H]
\caption{Regularisasi \textit{monotonic alignment} (contoh: KL terhadap topeng monotonic)}
\label{eq:mono}
\begin{equation}
\mathcal{L}_{\mathrm{mono}} \;=\;
\mathrm{KL}\!\left(A \;\middle\|\; M_{\mathrm{mono}}\right)
\;=\; \sum_{i,j} A_{ij}\,\log\!\frac{A_{ij}}{(M_{\mathrm{mono}})_{ij}},
\end{equation}
\end{eqfloat}

dengan $A$ matriks atensi dan $M_{\mathrm{mono}}$ topeng yang menegakkan urutan temporal (definisi spesifik mengikuti rancangan eksperimen).

% ---------- 2.6 Decoding + LM ----------
\subsection{Decoding dan Integrasi Model Bahasa}
Pada \textit{beam search}, skor akustik dari Whisper dapat dipadukan dengan LM eksternal (\textit{shallow fusion}) untuk menurunkan WER pada data terbatas (lihat bukti 2024 untuk bahasa under-represented).

\begin{eqfloat}[H]
\caption{Skor gabungan \textit{shallow fusion} untuk decoding}
\label{eq:sf}
\begin{equation}
\mathrm{score}(y\mid X) \;=\; 
\log P_\theta(y\mid X) \;+\; \lambda\,\log P_{\mathrm{LM}}(y) \;+\; \beta\,|y| ,
\end{equation}
\end{eqfloat}

dengan $\lambda,\beta$ hiperparameter dan $|y|$ penalti panjang.

% ---------- 2.7 Metrik evaluasi ----------
\subsection{Metrik Evaluasi}
\begin{eqfloat}[H]
\caption{Word Error Rate (WER)}
\label{eq:wer}
\begin{equation}
\mathrm{WER} \;=\; \frac{S + D + I}{N}\times 100\% ,
\end{equation}
\end{eqfloat}

\begin{eqfloat}[H]
\caption{Character Error Rate (CER)}
\label{eq:cer}
\begin{equation}
\mathrm{CER} \;=\; \frac{S_c + D_c + I_c}{N_c}\times 100\% .
\end{equation}
\end{eqfloat}

dengan $S,D,I$ adalah jumlah \textit{substitution}, \textit{deletion}, \textit{insertion} terhadap $N$ kata referensi; subskrip $_c$ untuk level karakter. CER sering lebih stabil pada bahasa dengan variasi ejaan atau morfologi kaya (praktik 2020–2025).

% ---------- 2.8 Keterkaitan ----------
\subsection{Keterkaitan dengan Penelitian Ini}
(1) Fondasi pra-latih skala besar \cite{baevski2020wav2vec2,conneau2021xlsr,chen2022wavlm,radford2023whisper} memotivasi adaptasi Whisper Small di skenario data terbatas.  
(2) Komparasi \textbf{full fine-tuning} vs \textbf{LoRA} akan mengevaluasi \textit{trade-off} akurasi–VRAM \cite{liu2024exploration}.  
(3) Jika baseline belum memadai, regularisasi \textbf{alignment monotonik} (Persamaan~\ref{eq:mono}) dapat dieksplorasi \cite{zhuang2025monotonic}.  
(4) \textbf{Shallow fusion} (Persamaan~\ref{eq:sf}) dan augmentasi sederhana akan dipertimbangkan untuk menekan WER pada Minangkabau.

